\chapter{Gốc Rễ Nhân Cách (The Roots of Character)}
\begin{center}
  \textit{\color{truyengold}``Nhân cách không được xây dựng trong những ngày bình an --- mà được tôi luyện trong cơn bão.''}\\[4pt]
  \textit{(Character is not built on peaceful days --- it is forged in the storm.)}\\[4pt]
  \small\color{truyenblue}--- Cổ Kim Đông Tây ---
\end{center}
\ngancach

\section{Mạnh Mẫu Ba Lần Dời Nhà}
\begin{truyen}{Mạnh Mẫu Ba Lần Dời Nhà}{Trung Quốc, thế kỷ 4 TCN}
\chuhoa{M}{ẹ} của Mạnh Tử lo lắng khi gia đình sống gần nghĩa địa: cậu bé chỉ biết chơi trò đào mộ và khóc thuê.\\[4pt]
\textit{(Mencius's mother was worried when the family lived near a cemetery: the boy only played at digging graves and wailing.)}

Bà dời nhà đến gần chợ; Mạnh Tử lại học đòi mặc cả và buôn bán.\\[4pt]
\textit{(She moved near a market; Mencius then imitated bargaining and trading.)}

Cuối cùng bà chuyển đến sát trường học; cậu bé bắt đầu học lễ nghĩa, học hành nghiêm túc.\\[4pt]
\textit{(Finally she moved near a school; the boy began learning etiquette and studying diligently.)}

Mạnh Mẫu dạy: ``Môi trường là người thầy đầu tiên và lớn nhất của một đứa trẻ.''\\[4pt]
\textit{(Mencius's mother taught: `Environment is a child's first and greatest teacher.')}

\end{truyen}
\begin{baihoc}
  Muốn con nên người, trước tiên hãy chọn môi trường cho con --- bởi ta trở thành những gì ta thường xuyên tiếp xúc.\\[4pt]
  \textit{(To raise a person of character, first choose the environment for them --- for we become what we are regularly exposed to.)}
\end{baihoc}

\section{Abe Lincoln Tự Học Dưới Ánh Lửa}
\begin{truyen}{Abe Lincoln Tự Học Dưới Ánh Lửa}{Abraham Lincoln --- Hoa Kỳ, 1820s}
\chuhoa{A}{braham} Lincoln lớn lên trong căn lều gỗ nghèo nàn; nhà không có đèn, không có trường, không có sách.\\[4pt]
\textit{(Abraham Lincoln grew up in a poor log cabin; no lamp, no school, no books.)}

Ông đi mượn sách từ nhà người hàng xóm cách xa hàng dặm, đọc dưới ánh lửa sưởi ấm vào buổi tối.\\[4pt]
\textit{(He walked miles to borrow books from neighbors, reading by firelight in the evenings.)}

Có lần ông mượn cuốn \textit{Tiểu sử Washington}; quyển sách bị ướt mưa, ông phải làm thuê ba ngày để đền bù.\\[4pt]
\textit{(Once he borrowed a biography of Washington; the book got rain-soaked and he worked three days to compensate for it.)}

Lincoln sau này nói: ``Thứ duy nhất tôi muốn biết là có thể tìm thấy trong sách.''\\[4pt]
\textit{(Lincoln later said: `The only thing I want to know can be found in books.')}

Từ đứa trẻ nghèo không được đến trường, ông trở thành vị tổng thống vĩ đại nhất nước Mỹ.\\[4pt]
\textit{(From a poor child with no schooling, he became the greatest president in American history.)}

\end{truyen}
\begin{baihoc}
  Hoàn cảnh quyết định điểm xuất phát --- nhưng ý chí và ham học mới quyết định điểm đến.\\[4pt]
  \textit{(Circumstances determine where you start --- but will and the love of learning determine where you end up.)}
\end{baihoc}

\section{Nguyễn Hiền Đỗ Trạng Nguyên Năm 13 Tuổi}
\begin{truyen}{Nguyễn Hiền Đỗ Trạng Nguyên Năm 13 Tuổi}{Nguyễn Hiền --- Việt Nam, thế kỷ 13}
\chuhoa{N}{guyễn} Hiền nhà nghèo đến nỗi không có tiền mua bút lông và giấy để học; ông dùng que tre viết trên đất, dùng lá cây làm tập.\\[4pt]
\textit{(Nguyen Hien's family was so poor he could not afford a brush and paper; he used bamboo sticks to write in soil and used leaves as notebooks.)}

Thầy giáo trong làng thương tình dạy miễn phí, nhưng Nguyễn Hiền không đủ ăn để đến lớp đều đặn.\\[4pt]
\textit{(The village teacher took pity and taught him for free, but Nguyen Hien often could not eat enough to attend regularly.)}

Thế mà năm 13 tuổi, ông dự thi và đỗ Trạng Nguyên --- người trẻ nhất trong lịch sử khoa cử Việt Nam.\\[4pt]
\textit{(Yet at 13, he sat the imperial exams and became the youngest Scholar-Champion in Vietnamese history.)}

Vua Trần hỏi tại sao không đến lạy kiến sớm hơn; Nguyễn Hiền thưa: ``Thần xin lỗi, thần không có giày.''\\[4pt]
\textit{(The Tran emperor asked why he had not come to pay respects sooner; Nguyen Hien replied: `I apologize, Your Majesty, I had no shoes.')}

\end{truyen}
\begin{baihoc}
  Thiếu thốn vật chất không thể cản được trí tuệ --- nếu trí tuệ đó được nuôi dưỡng bởi ý chí sắt đá.\\[4pt]
  \textit{(Material deprivation cannot stop the mind --- if that mind is nourished by an iron will.)}
\end{baihoc}

\section{Khổng Tử Học Không Biết Mỏi}
\begin{truyen}{Khổng Tử Học Không Biết Mỏi}{Khổng Tử --- Trung Quốc, thế kỷ 5 TCN}
\chuhoa{K}{hổng} Tử nói về bản thân: ``Ta không phải người sinh ra đã biết --- ta chỉ là người yêu học và tìm kiếm không ngừng.''\\[4pt]
\textit{(Confucius said of himself: `I was not born with knowledge --- I am simply one who loves learning and seeks endlessly.')}

Ông học nhạc từ một thầy già, luyện mãi một bản mà không chịu sang bản khác cho đến khi hiểu thấu cả tâm hồn người soạn nhạc.\\[4pt]
\textit{(He studied music from an elderly teacher, practicing one piece repeatedly, refusing to move on until he understood the very soul of the composer.)}

Bước sang tuổi 70, Khổng Tử vẫn nói: ``Mỗi sáng thức dậy tôi lại sợ rằng mình sẽ chết trước khi học xong những điều cần học.''\\[4pt]
\textit{(Into his seventies, Confucius still said: `Each morning I wake fearing I will die before I have learned what I need to learn.')}

\end{truyen}
\begin{baihoc}
  Người vĩ đại không phải là người giỏi nhất khi bắt đầu --- mà là người không bao giờ ngừng học.\\[4pt]
  \textit{(The great are not those who are best at the start --- but those who never stop learning.)}
\end{baihoc}
